\documentclass[border=0pt]{standalone}
% The shared preamble of every docs diagram (docs/diagrams/<name>.tex).
% dvisvgm's driver first, so patterns and opacity survive into the SVG.
\def\pgfsysdriver{pgfsys-dvisvgm.def}
\usepackage[siunitx,RPvoltages]{circuitikz}
\usepackage{xcolor}
% Latin Modern: outline fonts for every glyph a label uses, µ included.
\usepackage[T1]{fontenc}
\usepackage{lmodern}
\usetikzlibrary{arrows.meta,backgrounds,calc,patterns,positioning}
\tikzset{
  % A white card, so diagrams read in mdBook's dark themes too.
  every picture/.append style={show background rectangle,
    background rectangle/.style={fill=white,inner sep=4pt}},
  % Layout geometry: layers told apart by fill and hatching, never colour alone.
  layer a/.style={draw=black,line width=0.6pt,fill=gray!20},
  layer b/.style={draw=black,line width=0.6pt,fill=gray!45},
  layer c/.style={draw=black,line width=0.6pt,fill=white,postaction={pattern=north east lines}},
  layer d/.style={draw=black,line width=0.6pt,fill=white,postaction={pattern=crosshatch dots}},
  marker/.style={draw=black,dashed,line width=0.6pt},
  % A measured distance and its label.
  dim/.style={{Latex[length=2mm]}-{Latex[length=2mm]},line width=0.5pt},
  dim label/.style={font=\small,fill=white,inner sep=1pt},
  % Where a rule fails: thick, dashed and labelled, not only red.
  violation/.style={draw=red!75!black,line width=1.4pt,dash pattern=on 3pt off 2pt},
  violation label/.style={font=\small\bfseries,text=red!75!black},
  caption/.style={font=\small},
}

\begin{document}
\begin{circuitikz}
  % The example's two PMOS in one n-well; the well is tied to VDD.
  \draw (0,0) node[pmos,bulk,anchor=S] (P1) {};
  \draw (4,0) node[pmos,bulk,anchor=S] (P2) {};
  \draw (P1.S) -- ++(0,0.6) node[above] {VDD};
  \draw (P2.S) -- ++(0,0.6) node[above] {VDD2};
  \draw (P1.D) -- ++(0,-0.5) node[below] {OUT1};
  \draw (P2.D) -- ++(0,-0.5) node[below] {OUT2};
  \draw (P1.bulk) -- ++(0.6,0) coordinate (b1);
  \draw (P2.bulk) -- ++(0.6,0) coordinate (b2);
  \draw (b1) -- ++(0,-2.2) -| (b2);
  \draw ($(b1)+(0,-2.2)$) -- ++(-2.2,0) -- ++(0,3.2) -| (P1.S);
  \node[caption,anchor=north] at ($(b1)+(1,-2.2)$) {nwell (bulk) on VDD};
  \node[caption,anchor=east] at ($(P1.G)+(-0.2,0)$) {P1};
  \node[caption,anchor=east] at ($(P2.G)+(-0.2,0)$) {P2};
  \draw[violation] ($(P2)+(-0.9,-1)$) rectangle ($(P2)+(0.9,1)$);
  \node[violation label,anchor=west,align=left] at ($(P2)+(1.2,0)$) {no channel path\\ to its well's net};
\end{circuitikz}
\end{document}
